\appendix
% \noindent\textbf{Annotators.}\qquad
% All annotators are students from high level Chinese universities and we have tested their language skills to ensure that they can produce high quality annotations. During the annotation process, some of the annotators showed extremely outstanding linguistic skills, which also provided us with selection criteria for segmenting the dataset.\\
% \noindent\textbf{Pre-processing.}\qquad
% Before annotation, we manually cleaned the collected videos with two stages of clips. In the first stage, we cut all videos into content-independent clips, for example, we cut KT into individual performances. In the second stage, we first remove all videos that are unethical and contain sensitive information we cleaned the Humor videos to remove clips that do not meet ethical standards and have sensitive information, we do not operate on KT clips, only the score of each performance clip is recorded, and we subdivided the magic videos into clips to ensure that each clip includes only 1-3 coherent magic effects to avoid exceeding the model's understanding ability.\\
% \noindent\textbf{Tasks Annotation.}\qquad
% 1) We have developed detailed annotation documentation, and all annotators undergo rigorous training before performing the actual annotation. 
% 2) For the tasks in free-text, in Humor QA and MagicQA we suggested answers between 2-4 short sentences, in CreativeQA we would like a longer answer from the annotators. We asked the annotators to use concise and objective statements and avoid subjective adverbs and adjectives to ensure that the answers were in a similar language style. 
% 3) Annotators can question the counter-intuitive properties of video content or the generalizability of answers during the annotation process, and we will document such issues and address them in the next stage.
% 4) Cross-validation: We arranged another group of people who did not participate in the labeling to cross-validate the QA pair into three grades: completely wrong, ambiguous or too subjective, and no problem, and the statistical results can be seen in Figure \ref{}. because people's subjective understanding of humor, creativity and magic are bound to differ, a small number of disagreements is a normal result, we fixed the completely wrong QA pair and kept the rest.
% \noindent\textbf{Post-Processing.}\qquad
% After completing a round of annotation, we will check and filter the annotations. To ensure the quality of annotations, annotations with low quality or ambiguity will be returned to the task annotation stage for the next round of annotation. After all task annotations are completed with high quality, we will enter the data augmentation phase.


\section{More details of FunQA dataset} \label{appendix_a: FunQA more details}
\subsection{Dataset Construction Pipeline} \label{appendix_a:dataset con}

\begin{figure}[b]
    \centering
    \includegraphics[width=\linewidth]{figures_appendix/M4.pdf}
    \vspace{-0.6cm}
    \caption{\textbf{Example of Magic Method Task.} During the annotation process, we discovered that even as humans, it is difficult to fully understand the complete principles behind the implementation of magic tricks right from the beginning.}
    \vspace{-0.4cm}
    \label{fig:M4}
\end{figure}

\begin{figure}[!b]
    \centering
    \includegraphics[width=\linewidth]{figures_appendix/annotation_priciple.pdf}
    \vspace{-0.6cm}
    \caption{\textbf{FunQA Video Selection and Clip Principle.} We have a zero-tolerance policy regarding the inclusion of offensive content in our dataset. During the video sourcing process (video selection and video clipping stages), we ensure that such content is completely eliminated.}
    \vspace{-0.4cm}
    \label{fig:annotation-principle}
\end{figure}

\begin{figure}[!b]
    \centering
    \includegraphics[width=\linewidth]{figures_appendix/anno-error.pdf}
    \vspace{-0.6cm}
    \caption{\textbf{Examples of incorrect annotations and the suggested modifications.} Here we clarify the meanings of different cell colors. \textbf{Red}: Errors made by the annotator during the annotation process. \textbf{Yellow}: Types of annotation errors. \textbf{Green}: Reasons for the errors and suggested corrections.}
    \vspace{-0.4cm}
    \label{fig:annotation-error}
\end{figure}

\begin{figure}[!b]
    \centering
    \includegraphics[width=\linewidth]{figures_appendix/prompt-trans.pdf}
    \vspace{-0.6cm}
    \caption{\textbf{GPT prompt for Translation and augmentation.} Under our carefully designed prompt, GPT-3.5 can automatically generate English translations and five synonymous sentences from our high-quality Chinese annotations, thereby expanding FunQA dataset.}
    \vspace{-0.4cm}
    \label{fig:prompt-trans}
\end{figure}

\noindent\textbf{Pre-processing \& Qualification}\quad
For videos related to humor and magic, we downloaded them from different streaming platforms, mostly in the form of compilations. For Creative videos, we downloaded 26 episodes publicly available from Tokyo TV in Japan. We provided rigorous training to the annotators to ensure high-quality video clips in the final compilation. Annotators who successfully completed the Clip task according to the requirements are considered qualified and can proceed to the next stage of annotation.

\noindent\textbf{Training \& Annotation}\quad
We conducted systematic training for all annotators who passed the previous round of annotation, focusing on different tasks: \\
For the \textit{Counter-intuitive Timestamp Localization Task (H1, C1, and M1)}, the annotation format is a pair of numbers enclosed in square brackets, [xxxx, xxxx]. We asked the annotators to record the time intervals in which they felt pleasure (or amusement or shock) while watching the video. 

For the \textit{Task Detailed Description Task (H2, C2, and M2)}, we requested objective descriptions of what happened at [xxxx, xxxx], emphasizing a "what you see is what you get" approach. It is important to note that the annotations should only cover the selected time intervals and should not include subjective adverbs (such as vividly, vividly, or wildly). When describing characters or objects, be concise and add modifiers if there is ambiguity. An example of a poor label is "a man and a woman," while a relatively better example is "a man wearing a red hat and a woman wearing an apron." 

For the \textit{Counter-intuitiveness Reasoning Task (H3, C3, and M3)}, explain why the video is interesting in the context of the overall content. This part requires interpretive answers based on analysis, reasoning, and prior knowledge to explain why the video is counter-intuitive.

Specifically for the \textit{Magic Method Task (M4)}, we found that most annotators were not professional magicians, and even when watching instructional magic videos, it was challenging for them to provide complete and reasonable annotations for magic methods. Therefore, for M4, we only provided a small number of examples, which do not exist in the FunQA dataset. Figure \ref{fig:M4} shows an example of M4.

We conducted strict quality reviews in real-time for the annotations, including but not limited to checking if the annotations meet the requirements, if the selected videos contain sensitive content, and if the annotation texts contain sensitive content. The annotation principles (video selection and video clip, etc) can be seen in Figure~\ref{fig:annotation-principle}. Figure \ref{fig:annotation-error} also illustrates some errors encountered during the annotation process and the suggested modifications provided. After completing all the annotations, we conducted a consensus evaluation study to assess the objectivity of our annotations. We randomly selected 10\% of the videos and asked annotators to provide a consistency score (high consensus, medium consensus, low consensus) between their first and second annotations, considering all the previous annotations shown to them.

\noindent\textbf{Copyright \& License}\quad
We respect the copyright of each video. For content downloaded from streaming platforms, we obtained permission from the respective uploaders. For \textit{Kasou Taisho}, we submitted a fair use request to the official website of Tokyo TV in Japan.

We respect the personal identity information of everyone appearing in our videos and always strive to eliminate offensive content. By conducting strict reviews of the annotators and real-time quality checks during each annotation process, we ensure the absence of offensive content.
\subsection{Language Augmentation} \label{appendix_a: Augmentation}
For each QA pair, we first asked GPT-3.5 to translate the previous Chinese answers into English, and then let GPT-3.5 give five separate answers with the same meaning but different linguistic expressions. Afterward, we filtered out incorrect and incomplete generations. Figure \ref{fig:prompt-trans} shows the prompt we gave to GPT-3.5.

However, as mentioned in the main text in the limitation, the translation of GPT-3.5 will be ambiguous when the original text uses Chinese harmonics and some special Chinese meanings, and this problem mostly occurs in the Title Task, we screened such problems by hand and modified them, as shown in Figure \ref{fig:trans-err}.

\begin{figure}[htbp]
    \centering
    \includegraphics[width=\linewidth]{figures_appendix/tran_err.pdf}
    \vspace{-0.6cm}
    \caption{\textbf{Issues in translation caused by Chinese-English cultural differences.} In this example of translating and expanding a humorous video title, GPT-3.5 failed to understand the meaning of the original Chinese title, and we filtered out such data and made manual changes.}
    \vspace{-0.4cm}
    \label{fig:trans-err}
\end{figure}

\subsection{Data Statistics}\label{appendix_a: Dataset statics}
We collected raw videos from multiple video platforms, including short videos, long videos, and video clips, and the detailed data can be seen in Table \ref{T:raw_statics}. The word cloud shaped by FunQA logo of the all annotation word is shown in Figure \ref{fig:wc_funqa}.

% static table----------------------------
\begin{table}[h]
\caption{\textbf{Statistics of the FunQA raw data.} 
\label{T:raw_statics}}
\centering
\resizebox{1.0\linewidth}{!}{

\begin{tabular}{cccccc}
\hline
Raw Data          & Type                                & Source                   & Videos & Avg.len(s) & Total.len(h) \\ \hline
\multirow{3}{*}{Humor} & \multirow{2}{*}{Daily Life (Human)} & Bilibili                 & 351    & 182  & 15.8   \\
                       &                                     & TikTok                   & 1296   & 14      & 5.18    \\
                       & Nature (Animal)                     & Bilibili                 & 230    & 133  & 8.52    \\ \hline
Creative             & Performance                         & Kasou Taishou            & 26     & 6060  & 43.77   \\\hline
\multirow{2}{*}{MagicQA} & Close up Magic                      & \multirow{2}{*}{Youtube} & 765    & 96  & 20.40   \\ 
                       & Camera Magic                        &                          & 334    & 152  & 14.17   \\ \hline
FunQA                    & \textit{-}                          & -                        & 3002   & 129   & 107.87  \\ \hline\\
\end{tabular}

}
\end{table}
%static table ----------------------------------
\begin{figure}[htbp]
  \centering
  \includegraphics[width=0.5\textwidth]{figures/fig_word_cloud_funqa.png}
  \caption{\textbf{Word cloud of FunQA free-text answer.}}
  \label{fig:wc_funqa}
\end{figure}
